% Part: computability
% Chapter: computability-theory
% Section: halting-problem

\documentclass[../../../include/open-logic-section]{subfiles}

\begin{document}

\olfileid{cmp}{thy}{hlt}
\olsection{Masalah Mandheg}

Miturut konstruksine, fungsi komputabel parsial
universal~$\fn{Un}(e,x)$ ditegesi yen lan mung yen
komputasi fungsi mawa kode~$e$ ngasilake nilai
kanggo input~$x$. Lumrah yen kita takon apa
kahanan iki bisa diputusake. Kasunyatane, ora bisa.
Kanggo model komputasi mesin Turing, iki ateges
ora bisa diputusake sacara komputasional apa mesin
Turing tartamtu mandheg ing input tartamtu.
Mula teorema ing ngisor iki dikenal minangka
``ora bisane masalah mandheg diputusake''.
Kita bakal menehi rong bukti. Bukti kapisan
nerusake rembugan sadurunge, dene bukti
kapindho luwih langsung.

\begin{thm}
\ollabel{thm:halting-problem}
Tetepna
\[
h(e, x)  =
\begin{cases}
1 & \text{yen\/ $\fn{Un}(e, x)$ ditegesi} \\
0 & \text{yen ora.}
\end{cases}
\]
Mula $h$ ora komputabel.
\end{thm}

\begin{proof}
Upamane $h$ komputabel. Kita bakal nuduhake
yen iki ndadekake kita bisa netepake fungsi
komputabel universal. Tetepna
\[
\fn{Un'}(e,x) =
\begin{cases}
\fn{Un}(e,x) & \text{yen $h(e,x) = 1$} \\
0 & \text{yen ora.}
\end{cases}
\]
Nanging saiki $\fn{Un'}(e, x)$ fungsi total,
lan komputabel yen~$h$ komputabel. Umpamane,
kita bisa netepake $g$ nganggo rekursi primitif,
minangka
\begin{align*}
g(0, e, x) & \simeq 0 \\
g(y+1, e, x) & \simeq \fn{Un}(e,x);
\end{align*}
banjur
\[
\fn{Un'}(e,x) \simeq g(h(e,x),e,x).
\]
Amarga $\fn{Un'}(e,x)$ padha karo $\fn{Un}(e,x)$
ing saben papan kang ndadekake fungsi kapindho
iku ditegesi, $\fn{Un'}$ universal kanggo
fungsi komputabel parsial kang jebul total.
Nanging iki nalisir \olref[nou]{thm:no-univ}.
\end{proof}

\begin{proof}
Upamane $h(e,x)$ komputabel. Tetepna fungsi $g$
minangka
\[
g(x) =
\begin{cases}
  0                & \text{yen $h(x,x) = 0$} \\
  \fundefined & \text{yen ora.}
\end{cases}
\]
Fungsi $g$ komputabel parsial. Umpamane,
fungsi iki bisa ditegesi minangka
$\umin{y}{h(x,x) = 0}$. Mula, kanggo sawenehing~$e$,
$g(x) \simeq \fn{Un}(e, x)$ kanggo saben~$x$.
Apa $g$ ditegesi ing $e$? Yen iya, miturut
definisi~$g$, $h(e,e) = 0$ ($h$~mung bisa
duwé nilai~$0$ yen ditegesi). Miturut
definisi~$h$, iki ateges $\fn{Un}(e, e)$
ora ditegesi. Amarga kita nganggep
$g(x) \simeq \fn{Un}(e, x)$ kanggo saben~$x$,
mula $g(e)$ ora ditegesi, sawijining
kontradiksi. Kosok baline, yen $g(e)$ ora
ditegesi, mula $h(e,e) \neq 0$, saengga
$h(e,e) = 1$. Saka kene $\fn{Un}(e, e)$
ditegesi. Nanging amarga
$g(x) \simeq \fn{Un}(e, x)$, $g(e)$ uga
bakal ditegesi. Iki uga kontradiksi.
\end{proof}

\begin{tagblock}{TMs}
\begin{explain}
Argumen iki bisa dijlentrehake nganggo
mesin Turing. Upamane ana mesin Turing~$H$
kang nampa jlentrehan mesin Turing~$E$
lan input~$x$, banjur mutusake apa $E$
mandheg ing input~$x$. Kita banjur bisa
nggawe mesin Turing liya~$G$ kang nampa
siji input~$x$, nglakokake $H$ kanggo
mutusake apa mesin $M_x$ mawa indeks~$x$
mandheg ing input~$x$, banjur nindakake
kosok baline. Tegese, yen $H$ nglaporake
yen $M_x$ mandheg ing input~$x$, $G$
mlebu gelung tanpa wates; dene yen $H$
nglaporake yen $M_x$ ora mandheg ing
input~$x$, $G$ langsung mandheg. Apa
$G$ mandheg yen diwenehi indekse dhewe
minangka input? Argumen ing ndhuwur
nuduhake yen mesin iki mandheg yen lan
mung yen ora mandheg---kontradiksi.
Mula panganggep yen mesin Turing~$H$
kaya mangkono ana kudu salah.
\end{explain}
\end{tagblock}

\end{document}
